\section{Complete Experimental Design}
\label{app:experimental}

\subsection{Gate chain and provenance}

The experiment order is navigation capability $\rightarrow$ battery calibration $\rightarrow$ Oracle headroom $\rightarrow$ frozen-policy prediction $\rightarrow$ closed-loop decision $\rightarrow$ second-domain/filter replication $\rightarrow$ matched end-to-end CMDP comparison. A downstream gate cannot be reported when an upstream gate fails. Every reported cell must record configuration, code revision, checkpoint, raw artifact path, seed list, sample count, and evaluation script.

\subsection{Composition split}

Let \(\{\pi_i\}_{i=1}^P\) and \(\{\Pi_j\}_{j=1}^F\) be frozen policies and safety operators. Leave-one-pair-out evaluation withholds complete pairs while retaining separately seen components where possible. Within each held-out pair, states are stratified by an interface-risk score computed only from source executed \((x,u)\) samples and by true charger hitting horizon. The primary factorial analysis crosses pair seen/unseen, interface low/high extrapolation, and horizon short/long. Reporting the marginal pair effect without these controls is forbidden.

\subsection{Prediction and decision baselines}

\begin{table}[t]
\centering
\small
\caption{Baseline information contract. ``Query'' means access to the frozen target component, not target outcomes.}
\label{tab:baselines}
\begin{tabular}{lcccc}
\toprule
Method & $a^{\rm nom}$ & $u^{\rm exec}$ & query $\pi,\Pi$ & reliability \\
\midrule
Distance $\times$ rate & no & no & no & no \\
Receding-horizon planner & no & yes & yes & explicit model \\
Predictive-safety head & optional & yes & no & learned safety score \\
Direct MC / TD & optional & yes & no & ensemble option \\
PCM-Nominal & yes & no & $\pi$ & ensemble option \\
PCM-Executed & yes & yes & $\pi,\Pi$ & ensemble option \\
Executed-WM & no & yes & $\pi,\Pi$ & no \\
Executed-WM Ensemble & no & yes & $\pi,\Pi$ & ensemble \\
SIRP candidate & no & yes & $\pi,\Pi$ & support+horizon+ensemble \\
Oracle & realized & realized & yes & exact simulator clone \\
CMDP/DCRL track & task state & environment & no & method-specific \\
\bottomrule
\end{tabular}
\end{table}


The Oracle uses the simulator's realized charger Resource-to-Go under a cloned information state and the same target policy/filter, without exposing future task rewards to learned methods. The receding-horizon planner may query the same frozen plant/resource primitive but not realized future task rewards; its horizon, compute budget, and terminal approximation are reported. The predictive-safety head adapts the future-resource hierarchy of \citet{guo2020predictive} to the single-agent setting and receives no privileged target outcome. The direct switcher observes the same high-level state available to the ReturnManager but does not receive oracle future resource. CMDP methods receive a matched interaction budget and are evaluated under the same mission starts; CPO, CVPO, SDAC, and DCRL are included when their constraint semantics can be matched, and any non-equivalence is stated. Their collision semantics are reported separately from HOCBF-filtered modular methods.

\subsection{Metrics}

Prediction metrics are MAE, RMSE, signed bias, underestimation rate and magnitude, boundary-weighted underestimation, horizon-stratified error, accepted-set risk, coverage, and AURC. Mission metrics are energy exhaustion/stranding, tasks per battery cycle, tasks per simulated hour, successful charger return, unused resource at arrival, premature return, and the paired stranding--throughput frontier. CRPS, pinball losses, and conditional coverage are disabled unless a future stochastic-environment gate establishes a non-degenerate conditional resource law.

\subsection{Statistical reporting}

For a binary event with \(k\) occurrences in \(n\) cycles, report the Wilson interval at the declared confidence level. Before evaluation, choose \(n\) from the desired upper-bound or half-width precision at the declared stranding ceiling; the minimum of 100 cycles is only a floor, not evidence that a rare-event target has been resolved. For paired method contrasts, resample cycles with their shared task seeds and report percentile bootstrap intervals, number of resamples, and the complete seed set. The 5\% Oracle gate is reported together with its continuous effect estimate and a sensitivity analysis. No claim is formed from one navigation seed or fewer than the preregistered cycle count per frontier point.
